\section{Further boundary bounds and analytic predecessors}
\label{sec:predecessors}

Simpler constructions remain useful when they avoid certificate
dependencies, apply to particular input profiles, or resolve boundary
regimes. We begin with a greedy construction for inputs without heavy
modes.

\begin{lemma}[Distinct greedy blocks for light inputs]\label{lem:all-light}
Suppose $m_i\le E$ for all modes, where $E>0$. For every $1\le k<n$,
at most $k$ distinct blocks reach
\begin{equation}\label{eq:all-light-reach}
 \min\left\{T,\frac{k(n-k+1)}{n-k}E\right\}
\end{equation}
with full error at most $E$ on that prefix.
\end{lemma}
\begin{proof}
Put $t_0=0$ and $t_j=j(n-j+1)E/(n-j)$ for $1\le j\le k$.
These times are strictly increasing: the function
$j(n-j+1)/(n-j)=j+j/(n-j)$ increases on the integers below $n$.
Suppose $j-1$ distinct modes have reached $t_{j-1}<T$. At
$u=\min\{T,t_j\}$ their allocations sum to at most $(j-1)E$.
One of the $n-j+1$ unused modes therefore satisfies
\[
 A_i(u)\ge\frac{u-(j-1)E}{n-j+1}.
\]
The identity
\[
 (n-j)t_j=(n-j+1)t_{j-1}+(n-2j+2)E
\]
and $u\le t_j$ imply
\[
 \frac{u-(j-1)E}{n-j+1}\ge u-t_{j-1}-E.
\]
Activate this mode from $t_{j-1}$ to $u$. It was previously unused,
so its endpoint discrepancy is $u-t_{j-1}-A_i(u)\le E$; during the
block this discrepancy is nondecreasing. Earlier selected modes receive
no additional service and preserve their one-sided bound. All positive
discrepancies are bounded by their terminal masses, hence by $E$.
Stop at $T$, or continue the induction.
\end{proof}

\begin{corollary}[Boundary consequences and an equal-mass band]
\label{cor:light-boundaries}
For every $k\ge1$ and $2\le n\le k$,
\begin{equation}\label{eq:heavy-boundary}
 F_{n,k-1}(T)\le\frac{T}{k+1}.
\end{equation}
For $k+1\le n\le2k$, the exact value is $F_{n,k-1}(T)=T/(k+1)$.
Furthermore, every input with terminal masses $T/n$ has instance optimum
$\OPT_{k-1}(A)=T/n$ throughout the band
\begin{equation}\label{eq:equal-mass-band}
 1\le k<n,\qquad (n-k)^2\le k,
 \qquad G^=_{n,k}(T)=\frac Tn.
\end{equation}
\end{corollary}
\begin{proof}
At $E=T/(k+1)$ and $n\le k$, some mass is at least
$T/n>E$. Theorem~\ref{thm:heavy} applies directly with $s=k-1$.
The all-light case is impossible here and is not needed; no matching
lower bound is claimed for~\eqref{eq:heavy-boundary}.
For $k<n\le2k$, either the same heavy argument applies or
Lemma~\ref{lem:all-light} reaches $T$, since
$k(n-k+1)/(n-k)\ge k+1$ is equivalent to $n\le2k$.
The $k+1$-pure-mode witness gives equality. For $k=1$ this includes
$n=2$; for $k=2$ it gives $n=3,4$. Some of these endpoints go beyond
the elementary plateau interval, though Theorem~\ref{thm:small-full}
already covers them.

For equal terminal masses, instead set $E=T/n$. Writing $d=n-k>0$,
the reach condition becomes $k(d+1)\ge nd=(k+d)d$, or $k\ge d^2$.
Lemma~\ref{lem:all-light} therefore gives full error at most $T/n$.
Every $k$-block schedule omits a mode, whose terminal mass is $T/n$;
this proves the lower bound for every input in this class. The band
includes $n=k+2$ for all $k\ge4$, $n=k+3$ for all $k\ge9$, and in
general all fixed deficits $d$ once $k\ge d^2$. The statement does
not assert a characterization outside this band.
\end{proof}

\begin{proposition}[The dimension-free limit]\label{prop:dimension-free}
For every fixed $k\ge1$,
\[
 F_{n,k-1}(T)<T/k\quad\hbox{for every finite }n\ge2,
 \qquad \sup_{n\ge2}F_{n,k-1}(T)=T/k.
\]
\end{proposition}
\begin{proof}
If $n\le k$, use~\eqref{eq:heavy-boundary}. If $n>k$, then
\[
 \frac1k-C_{n,k}
 =\frac{(k+1)(n-k)}{nk(2n-k-1)}>0,
\]
and both entries of the full upper bound~\eqref{eq:general-full} are
strictly below $1/k$. The uniform lower coefficient tends to $1/k$,
proving the supremum assertion. The leading dimension-free estimate also
has a direct classical rounding proof: divide the horizon into $k$ equal
slots and use the network in Lemma~\ref{lem:prefix-repeat}, without its
objective requiring a repetition. Integral prefix counts lie between
floor and ceiling allocations. Each endpoint error is strictly smaller
than a slot length, and Lemma~\ref{lem:endpoint} preserves that error
inside each slot. This explains why the leading limit is already a
coarse-slot rounding consequence; the finite-mode corrections require
the sharper constructions above.
\end{proof}

The classical unrestricted CIA estimate also gives a useful explicit
comparison for every $n\ge2$ and $k\ge1$:
\begin{equation}\label{eq:classical-small-mode}
 F_{n,k-1}(T)\le\frac{2n-3}{2n-2}\frac{T}{k}.
\end{equation}
Indeed, apply \citet[Corollary~1, p.~669]{zeile2021dwell} to the cell averages
on $k$ equal cells. Every integer choice on these cells has at most $k-1$
switches, and cell averaging preserves its error for the original measurable
input. The source's sharp unrestricted-grid constant does not imply that
\eqref{eq:classical-small-mode} is sharp with this hard switch budget. For
fixed small $n$ and sufficiently large $k$, it improves $T/(k+1)$; it also
makes the classical origin of the leading dimension-free estimate explicit.

\begin{proposition}[A certificate-free four-block upper bound]
\label{prop:analytic-four-upper}
For $n\ge5$, define
\begin{equation}\label{eq:analytic-four-upper}
 U_n=\frac{n^4-7n^3+21n^2-29n+15}
 {n(2n-3)(2n^2-6n+5)}.
\end{equation}
Every input admits four distinct blocks with one-sided error at most
$TU_n$, and $F_{n,3}(T)\le T\max\{1/5,U_n\}$.
This proves the exact plateau $F_{n,3}(T)=T/5$ for $5\le n\le11$
without the four-block certificates.
\end{proposition}
\begin{proof}
Use the analytic three-block seed $L_{n-1,3}$ and apply
Lemma~\ref{lem:mode-removal} once. Substitution gives
$f_n(L_{n-1,3})=U_n$. At $n=5$ the seed is below $1/(n-1)$,
so the minimum-mass branch supplies the one-sided conclusion there;
for $n\ge6$ the maximum-mass branch applies. The denominator in
\eqref{eq:analytic-four-upper} is positive for $n\ge5$.
The full bound follows from the heavy-mode reduction.

The numerator of $U_n-1/5$, after clearing its positive denominator, is
\[
 p(n)=n^4-17n^3+77n^2-130n+75.
\]
At $n=5,6,7,8,9,10,11$ its values are respectively
$-150,-309,-492,-645,-690,-525,-24$. At $n=12+h$ it is
\[
 h^4+31h^3+329h^2+1286h+963>0\qquad(h\ge0).
\]
Thus exactly the stated integer interval satisfies $U_n\le1/5$.
The omitted-mode witness gives the matching lower bound throughout it.
\end{proof}

The exact value~\eqref{eq:three-switch-exact} is stronger, but this bound
avoids the four-block certificates.
For a quantitative comparison,
\begin{equation}\label{eq:analytic-four-gap}
 U_n-L_{n,4}=
 \frac{6n^4-20n^3+21n^2-7n+1}
 {(2n-3)(2n-1)(2n^2-6n+5)(2n^2-2n+1)}>0.
\end{equation}
Positivity follows by writing the numerator at $n=5+h$ as
$6h^4+100h^3+621h^2+1703h+1741$. Power-series division gives
\[
 U_n=\frac14-\frac5{8n}+\frac{11}{16n^2}+O(n^{-3}),\qquad
 U_n-L_{n,4}=\frac3{8n^2}+O(n^{-3}).
\]
For $n\ge12$ the latter is precisely the gap above the exact full
minimax coefficient; on the plateau both full bounds equal $1/5$.

\begin{lemma}[A sufficient condition using the third largest mass]
\label{lem:third-mass}
Let $n\ge3$, let $m_{(3)}$ be the third largest terminal mass, and put
$r=n/(n-1)$. If
\begin{equation}\label{eq:third-mass}
 \max_i m_i\le E,\qquad
 E\ge\frac{T-m_{(3)}}{1+r+r^2},
\end{equation}
then at most three distinct blocks suffice for full error at most $E$.
\end{lemma}
\begin{proof}
Set $L=T-m_{(3)}-E$. If $L\le0$, any mode among the three largest
masses can be used throughout: its one-sided error is at most
$T-m_{(3)}\le E$, and all positive discrepancies are bounded by the
masses. If $L>0$, then $L<T$, and~\eqref{eq:third-mass} is exactly
$L\le(r+r^2)E=B_2(n,E)$. Lemma~\ref{lem:reach-two} supplies at most
two distinct blocks on $[0,L]$. At least one of the three largest-mass
modes is unused. Append it on $[L,T]$; its final one-sided discrepancy
is at most $T-L-m_{(3)}=E$ and is nondecreasing during its block.
Other one-sided discrepancies can only decrease afterwards. The mass
hypothesis bounds every positive discrepancy.
\end{proof}

This condition can improve on an unrestricted minimax guarantee for
individual inputs. For example, at $n=8$ take terminal masses $1/6$
for three modes and $1/10$ for the other five, on horizon one. Then
$E=245/1014$ satisfies~\eqref{eq:third-mass}, is larger than all the
masses, and is strictly smaller than $L_{8,3}=343/1352$.
Thus every temporal profile with these masses enjoys the stronger
instance guarantee, although the exact unrestricted value is the latter
coefficient.

